%% ma: L11-B27-0020
%% lop: 11
%% chuong: 7
%% bai: 27
%% dang: 11.27.16
%% loai: TN4
%% muc: VD
%% dap_an: "B"
%% nguon: "(NBV)-ĐỀ SỐ 18-MỖI NGÀY 1 ĐỀ THI 2025.pdf (D18, MỖI NGÀY 1 ĐỀ - Đề số 18), Phần I Câu 11"
%% kien_thuc: "Bài 27 (thể tích khối chóp, khối lăng trụ), Bài 26 (khoảng cách từ điểm đến mặt phẳng)"
%% trang_thai: chinh-thuc
%% ghi_chu: "Đáp án gốc B đúng (V = 1/2 . 15 . 6 = 45, giải độc lập: V(C.ABB'A') = 1/3 . 15 . 6 = 30 = 2/3 V). Hình dựng lại lăng trụ xiên, nét khuất AC. Đã duyệt 19/10/2026."
\begin{ex}
Cho lăng trụ tam giác $ABC.A'B'C'$. Biết diện tích mặt bên $(ABB'A')$ bằng $15$, khoảng cách từ $C$ đến mặt phẳng $(ABB'A')$ bằng $6$. Thể tích của khối lăng trụ $ABC.A'B'C'$ bằng bao nhiêu? (tham khảo hình vẽ bên dưới)
\begin{center}
\begin{tikzpicture}[scale=.9]
  \coordinate (A) at (0,0);\coordinate (B) at (1.4,-.9);\coordinate (C) at (3.2,0);
  \coordinate (A') at (.5,2.6);\coordinate (B') at (1.9,1.7);\coordinate (C') at (3.7,2.6);
  \draw (A)--(B)--(C) (A')--(B')--(C')--cycle (A)--(A') (B)--(B') (C)--(C');
  \draw[khuat] (A)--(C);
  \foreach \p/\d in {A/left,B/below,C/right,A'/above left,B'/below right,C'/above right}
    \fill (\p) circle (1pt) node[\d,font=\footnotesize] {$\p$};
\end{tikzpicture}
\end{center}
\choice{$60$}{\True $45$}{$90$}{$30$}
\loigiai{Thể tích khối chóp $C.ABB'A'$ là $V_1=\dfrac13\cdot15\cdot6=30$.

Khối lăng trụ tam giác $ABC.A'B'C'$ gồm khối chóp $C.ABB'A'$ và khối chóp $C.A'B'C'$, mà thể tích khối chóp $C.A'B'C'$ bằng $\dfrac13V$ ($V$ là thể tích lăng trụ). Do đó $V_1=\dfrac23V$, suy ra $V=\dfrac32V_1=\dfrac32\cdot30=45$. Chọn B.}
\end{ex}
